% Part: computability
% Chapter: recursive-functions
% Section: pr-relations

\documentclass[../../../include/open-logic-section]{subfiles}

\begin{document}

% SEGMENT OLP-0217-S01

\olfileid{cmp}{rec}{prr}
\olsection{முதன்மை மீள்வரையறை உறவுகள்}


\begin{defn}
ஓர் உறவு $R(\vec x)$ இன் சிறப்பியல்புச் சார்பு
\[
\Char{R}(\vec x) = \left\{
  \begin{array}{ll}
  1 & \mbox{if $R(\vec x)$} \\
  0 & \mbox{இல்லையெனில்}
  \end{array}
\right.
\]
முதன்மை மீள்வரையறைச் சார்பாக இருந்தால், அந்த உறவு முதன்மை
மீள்வரையறை உறவு எனப்படும்.
\end{defn}

வேறுவிதமாகக் கூறினால், முதன்மை மீள்வரையறை உறவு $R(\vec x)$
ஒன்றைக் குறிப்பிடும்போது, $\Char{R}(\vec x) = 1$ என்ற வடிவிலான
உறவையே குறிப்பிடுகிறோம்; இங்கே $\Char{R}$ ஒவ்வோர் உள்ளீட்டுக்கும்
1 அல்லது 0 ஐத் திருப்பித்தரும் ஒரு முதன்மை மீள்வரையறைச் சார்பு.
எடுத்துக்காட்டாக, $x = 0$ எனில், எனில் மட்டுமே, பொருந்தும் உறவு
$\fn{IsZero}(x)$ ஆனது, முதன்மை மீள்வரையறையைப் பயன்படுத்திப்
பின்வருமாறு வரையறுக்கப்படும் சார்பு $\Char{\fn{IsZero}}$ க்கு ஒத்தது:
\begin{align*}
\Char{\fn{IsZero}}(0) & = 1,\\
\Char{\fn{IsZero}}(x+1) & = 0.
\end{align*}

உறவுகளைப் பிற முதன்மை மீள்வரையறைச் சார்புகளுடன் சேர்க்கலாம் என்பது
தெளிவாக இருக்க வேண்டும். எனவே பின்வருவனவும் முதன்மை
மீள்வரையறை உறவுகள்:
\begin{enumerate}
\item $\fn{IsZero}(\left|x - y\right|)$ ஆல் வரையறுக்கப்படும் சம
உறவு, $x = y$
\item $\fn{IsZero}(x \tsub y)$ ஆல் வரையறுக்கப்படும் குறைவு-அல்லது-சம
உறவு, $x \leq y$
\end{enumerate}


\begin{prop}
  முதன்மை மீள்வரையறை உறவுகளின் கணம் பூலியச் செயல்களின் கீழ்
  மூடியது; அதாவது, $P(\vec x)$ மற்றும் $Q(\vec x)$ முதன்மை
  மீள்வரையறை உறவுகள் எனில், பின்வருவனவும் அவ்வாறே அமையும்:
  \begin{enumerate}
  \item $\lnot P(\vec x)$
  \item $P(\vec x) \land Q(\vec x)$
  \item $P(\vec x) \lor Q(\vec x)$
  \item $P(\vec x) \lif Q(\vec x)$
  \end{enumerate}
\end{prop}

\begin{proof}
  $P(\vec x)$ மற்றும் $Q(\vec x)$ முதன்மை மீள்வரையறை உறவுகள்,
  அதாவது அவற்றின் சிறப்பியல்புச் சார்புகள் $\Char{P}$ மற்றும்
  $\Char{Q}$ அவ்வாறானவை, என்று கொள்க. $\lnot P(\vec x)$
  முதலியவற்றின் சிறப்பியல்புச் சார்புகளும் முதன்மை மீள்வரையறைச்
  சார்புகள் என்று காட்ட வேண்டும்.
  \[
  \Char{\lnot P}(\vec x) = \begin{cases}
    0 & \text{$\Char{P}(\vec x) = 1$ எனில்}\\
    1 & \text{இல்லையெனில்}
  \end{cases}
  \]
  $\Char{\lnot P}(\vec x)$ ஐ $1 \tsub \Char{P}(\vec x)$ என
  வரையறுக்கலாம்.
  \[
  \Char{P \land Q}(\vec x) = \begin{cases}
    1 & \text{$\Char{P}(\vec x) = \Char{Q}(\vec x) = 1$ எனில்}\\
    0 & \text{இல்லையெனில்}
  \end{cases}
  \]
  $\Char{P \land Q}(\vec x)$ ஐ $\Char{P}(\vec x) \cdot
  \Char{Q}(\vec x)$ எனவோ, $\fn{min}(\Char{P}(\vec x), \Char{Q}(\vec
  x))$ எனவோ வரையறுக்கலாம். அதேபோல்,
  \begin{align*}
    \Char{P \lor Q}(\vec x) & = \fn{max}(\Char{P}(\vec x), \Char{Q}(\vec x)) \text{ மற்றும்}\\
    \Char{P \lif Q}(\vec x) & = \fn{max}(1 \tsub
  \Char{P}(\vec x), \Char{Q}(\vec x)).
  \end{align*}
\end{proof}

\begin{prop}
  முதன்மை மீள்வரையறை உறவுகளின் கணம் வரம்பிட்ட அளவையிடலின் கீழ்
  மூடியது; அதாவது, $R(\vec x, z)$ ஒரு முதன்மை மீள்வரையறை உறவு
  எனில், பின்வரும் உறவுகளும் அவ்வாறே அமையும்:
  \begin{align*}
    & \bforall{z < y}{R(\vec x, z)} \text{ மற்றும்}\\
    & \bexists{z < y}{R(\vec x, z)}.
  \end{align*}
  ஒவ்வொரு~$y$ ஐவிடக் குறைந்த~$z$ க்கும் $R(\vec x, z)$
  பொருந்தினால், எனில் மட்டுமே, $\bforall{z < y}{R(\vec x, z)}$
  ஆனது $\vec x$ மற்றும் $y$ க்குப் பொருந்தும்; இதேபோல்
  $\bexists{z < y}{R(\vec x, z)}$ க்கும்.
\end{prop}

\begin{proof}
  மரபின்படி, $\bforall{z < 0}{R(\vec x, z)}$ மெய் என்று கொள்கிறோம்
  (~$0$ ஐவிடக் குறைந்த $z$ எதுவும் இல்லை என்ற எளிய காரணத்தால்);
  மேலும் $\bexists{z < 0}{R(\vec x, z)}$ பொய் என்று கொள்கிறோம்.
  வரம்பிட்ட அனைத்தளவையீட்டுக்குறி ஓர் இறுதிக்குட்பட்ட பெருக்குத்தொகை
  அல்லது மீண்டும் மீண்டும் எடுக்கப்படும் மீச்சிறு மதிப்பு போலவே
  செயல்படுகிறது. அதாவது, $P(\vec x, y) \defiff \bforall{z <
  y}{R(\vec x, z)}$ எனில், $\Char{P}(\vec x, y)$ ஐப் பின்வருமாறு
  வரையறுக்கலாம்:
  \begin{align*}
    \Char{P}(\vec x, 0) & = 1\\
    \Char{P}(\vec x, y+1) & =
    \fn{min}(\Char{P}(\vec x, y), \Char{R}(\vec x, y)).
  \end{align*}
  வரம்பிட்ட உளதளவையிடலை ஒத்தவாறு $\fn{max}$ ஐப் பயன்படுத்தி
  வரையறுக்கலாம். மாறாக, சமானம் $\bexists{z < y}{R(\vec x, z)}
  \liff \lnot \bforall{z < y}{\lnot R(\vec x, z)}$ ஐப் பயன்படுத்தி,
  வரம்பிட்ட அனைத்தளவையிடலிலிருந்தும் அதை வரையறுக்கலாம்.
  எடுத்துக்காட்டாக, $\bexists{x \leq y}{\dots x\dots}$ என்ற
  வடிவிலான வரம்பிட்ட அளவையீட்டுக்குறி
  $\bexists{x < y+1}{\dots x \dots}$ குச் சமானமானது என்பதைக் கவனிக்க.
\end{proof}

\begin{prob}
  மூவிட உறவு $x \equiv y \mod n$ (மாடுலோ~$n$ ஒப்புமை) ஒரு
  முதன்மை மீள்வரையறை உறவு என்று காட்டுக.
\end{prob}

பின்வருமாறு வரையறுக்கப்படும் நிபந்தனைச் சார்பு
$\fn{cond}(x,y,z)$ மற்றொரு பயனுள்ள முதன்மை மீள்வரையறைச் சார்பு:
\begin{align*}
  \fn{cond}(x,y,z) & = \begin{cases}
  y & \text{$x = 0$ எனில்} \\
  z & \text{இல்லையெனில்}.
\end{cases}
\intertext{இது பின்வருமாறு மீள்வரையறை முறையில் வரையறுக்கப்படுகிறது:}
\fn{cond}(0,y,z) & = y,\\
\fn{cond}(x+1,y,z) & = z.
\end{align*}
முதன்மை மீள்வரையறை உறவுகளிலிருந்து நேர்வுவாரியாக முதன்மை
மீள்வரையறைச் சார்புகளை வரையறுப்பதை நியாயப்படுத்த இதைப் பயன்படுத்தலாம்:

\begin{prop}
$g_0(\vec x)$, \dots,~$g_m(\vec x)$ ஆகியவை முதன்மை
மீள்வரையறைச் சார்புகள் என்றும், $R_0(\vec x)$, \dots,
$R_{m-1}(\vec x)$ ஆகியவை முதன்மை மீள்வரையறை உறவுகள் என்றும்
இருந்தால், பின்வருமாறு வரையறுக்கப்படும் சார்பு $f$:
\[
f(\vec x) = \begin{cases}
    g_0(\vec x) & \text{$R_0(\vec{x})$ எனில்} \\
    g_1(\vec x) & \text{$R_1(\vec{x})$ பொருந்தி $R_0(\vec{x})$ பொருந்தாவிடில்} \\
    \vdots & \\
    g_{m-1}(\vec x) & \text{$R_{m-1}(\vec{x})$ பொருந்தி முந்தையவை
      எதுவும் பொருந்தாவிடில்}
    \\
    g_m(\vec x) & \mbox{இல்லையெனில்}
\end{cases}
\]
ஒரு முதன்மை மீள்வரையறைச் சார்பாகவும் அமையும்.
\end{prop}

\begin{proof}
  $m = 1$ எனில், இது பின்வருமாறு வரையறுக்கப்படும் சார்பு மட்டுமே:
  \[
  f(\vec x) = \fn{cond}(\Char{\lnot R_0}(\vec x),g_0(\vec x),g_1(\vec
  x)).
  \]
  $m$ ஆனது $1$ ஐவிடப் பெரியதாக இருந்தால், இவ்வடிவிலான
  வரையறைகளைச் சேர்த்தால் போதும்.
\end{proof}
\end{document}
